\documentclass[12pt,a4paper]{article}
\usepackage[UTF8]{ctex}
\usepackage{amsmath,amssymb,amsthm}
\usepackage{geometry}

\geometry{margin=2.5cm}

\title{混合运动-力控制器详解}
\author{第11章补充说明}
\date{}

\begin{document}

\maketitle

\section{问题提出}

本文档通俗地讲解"混合运动-力控制器"这一节，并详细论证其中的所有公式。我们将从先验知识开始，逐步推导出混合控制器的完整形式。

\section{先验知识回顾}

\subsection{任务空间动力学方程}

\textbf{无约束情况下的任务空间动力学}：
\begin{equation}
F = \Lambda(\theta)\dot{V} + \eta(\theta, V) \tag{11.55}
\end{equation}

其中：
\begin{itemize}
\item $F \in \mathbb{R}^6$：任务空间wrench（力和力矩）
\item $\Lambda(\theta) \in \mathbb{R}^{6 \times 6}$：任务空间质量矩阵
\item $\dot{V} \in \mathbb{R}^6$：末端执行器加速度（twist的导数）
\item $\eta(\theta, V) \in \mathbb{R}^6$：任务空间的科里奥利力、向心力和重力项
\end{itemize}

\textbf{推导过程}（简要）：
\begin{enumerate}
\item 从关节空间动力学：$\tau = M(\theta)\ddot{\theta} + h(\theta, \dot{\theta})$
\item 使用速度关系：$V = J(\theta)\dot{\theta}$
\item 使用加速度关系：$\dot{V} = \dot{J}(\theta)\dot{\theta} + J(\theta)\ddot{\theta}$
\item 应用虚功原理：$F = J^{-T}\tau$
\item 定义：$\Lambda(\theta) = J^{-T}M(\theta)J^{-1}$，$\eta(\theta, V) = J^{-T}h(\theta, J^{-1}V) - \Lambda\dot{J}J^{-1}V$
\end{enumerate}

详细推导见文档：\texttt{task\_space\_dynamics\_derivation.tex}

\subsection{Pfaffian约束}

\textbf{环境约束}：
\begin{equation}
A(\theta)V = 0 \tag{11.54}
\end{equation}

其中：
\begin{itemize}
\item $A(\theta) \in \mathbb{R}^{k \times 6}$：约束矩阵（$k$个约束）
\item $V \in \mathbb{R}^6$：末端执行器twist（速度）
\item $k$：约束的数量
\end{itemize}

\textbf{物理意义}：
\begin{itemize}
\item 环境在$k$个方向上限制了末端执行器的运动
\item 这些约束限制了速度$V$，使其必须满足$A(\theta)V = 0$
\end{itemize}

\section{约束动力学的推导}

\subsection{问题：如何考虑约束？}

\textbf{无约束动力学}：
\begin{equation}
F = \Lambda(\theta)\dot{V} + \eta(\theta, V)
\end{equation}

\textbf{问题}：当有约束$A(\theta)V = 0$时，动力学方程如何修改？

\textbf{答案}：使用拉格朗日乘数法！

\subsection{拉格朗日乘数法}

\textbf{基本思想}：
\begin{itemize}
\item 约束$A(\theta)V = 0$限制了系统的运动
\item 约束会产生约束力$F_{\text{constraint}}$
\item 约束力与约束方向垂直（或平行，取决于定义）
\end{itemize}

\textbf{约束动力学}：
\begin{equation}
F = \Lambda(\theta)\dot{V} + \eta(\theta, V) + F_{\text{constraint}} \tag{11.56}
\end{equation}

\textbf{约束力的形式}：
\begin{itemize}
\item 约束力必须与约束方向相关
\item 使用拉格朗日乘数：$F_{\text{constraint}} = A^T(\theta)\lambda$
\item 其中$\lambda \in \mathbb{R}^k$是拉格朗日乘数向量
\end{itemize}

\textbf{最终形式}：
\begin{equation}
F = \Lambda(\theta)\dot{V} + \eta(\theta, V) + A^T(\theta)\lambda \tag{11.56}
\end{equation}

\textbf{为什么是$A^T(\theta)$？}
\begin{itemize}
\item 约束$A(\theta)V = 0$定义了约束子空间
\item 约束力$F_{\text{constraint}}$必须在约束子空间的正交补空间中
\item $A^T(\theta)$的列空间就是约束子空间的正交补空间
\item 因此$F_{\text{constraint}} = A^T(\theta)\lambda$是合理的
\end{itemize}

\subsection{约束的时间导数}

\textbf{约束必须在所有时间满足}：
\begin{equation}
A(\theta)V = 0 \quad \text{（对所有时间$t$）}
\end{equation}

\textbf{对时间求导}：
\begin{align}
\frac{d}{dt}[A(\theta)V] &= 0 \\
\dot{A}(\theta)V + A(\theta)\dot{V} &= 0
\end{align}

\textbf{结果}：
\begin{equation}
A(\theta)\dot{V} + \dot{A}(\theta)V = 0 \tag{11.57}
\end{equation}

\textbf{物理意义}：
\begin{itemize}
\item 如果速度满足约束，加速度也必须满足约束
\item 这确保了约束在所有时间都满足
\end{itemize}

\section{拉格朗日乘数的求解}

\subsection{问题：如何求解$\lambda$？}

\textbf{已知}：
\begin{itemize}
\item 约束动力学：$F = \Lambda(\theta)\dot{V} + \eta(\theta, V) + A^T(\theta)\lambda$
\item 约束的时间导数：$A(\theta)\dot{V} + \dot{A}(\theta)V = 0$
\end{itemize}

\textbf{目标}：求解$\lambda$

\subsection{求解步骤}

\textbf{步骤1：从约束动力学求解$\dot{V}$}

从方程(11.56)：
\begin{equation}
F = \Lambda(\theta)\dot{V} + \eta(\theta, V) + A^T(\theta)\lambda
\end{equation}

求解$\dot{V}$：
\begin{align}
\Lambda(\theta)\dot{V} &= F - \eta(\theta, V) - A^T(\theta)\lambda \\
\dot{V} &= \Lambda^{-1}(\theta)[F - \eta(\theta, V) - A^T(\theta)\lambda] \tag{1}
\end{align}

\textbf{步骤2：代入约束的时间导数}

将方程(1)代入约束的时间导数(11.57)：
\begin{align}
A(\theta)\dot{V} + \dot{A}(\theta)V &= 0 \\
A(\theta)\Lambda^{-1}(\theta)[F - \eta(\theta, V) - A^T(\theta)\lambda] + \dot{A}(\theta)V &= 0
\end{align}

\textbf{步骤3：展开并整理}

展开：
\begin{align}
A(\theta)\Lambda^{-1}(\theta)F - A(\theta)\Lambda^{-1}(\theta)\eta(\theta, V) - A(\theta)\Lambda^{-1}(\theta)A^T(\theta)\lambda + \dot{A}(\theta)V &= 0
\end{align}

整理关于$\lambda$的项：
\begin{align}
A(\theta)\Lambda^{-1}(\theta)A^T(\theta)\lambda &= A(\theta)\Lambda^{-1}(\theta)[F - \eta(\theta, V)] + \dot{A}(\theta)V
\end{align}

\textbf{步骤4：求解$\lambda$}

假设$A(\theta)\Lambda^{-1}(\theta)A^T(\theta)$可逆（通常成立，因为约束是独立的），我们有：
\begin{equation}
\lambda = [A(\theta)\Lambda^{-1}(\theta)A^T(\theta)]^{-1}\left\{A(\theta)\Lambda^{-1}(\theta)[F - \eta(\theta, V)] + \dot{A}(\theta)V\right\} \tag{2}
\end{equation}

\textbf{步骤5：简化}

从约束$A(\theta)V = 0$，我们知道$\dot{A}(\theta)V = -A(\theta)\dot{V}$（从约束的时间导数）。

但更直接地，我们可以使用：
\begin{equation}
\dot{A}(\theta)V = -A(\theta)\dot{V}
\end{equation}

代入方程(2)：
\begin{align}
\lambda &= [A(\theta)\Lambda^{-1}(\theta)A^T(\theta)]^{-1}\left\{A(\theta)\Lambda^{-1}(\theta)[F - \eta(\theta, V)] - A(\theta)\dot{V}\right\} \\
        &= [A(\theta)\Lambda^{-1}(\theta)A^T(\theta)]^{-1}\left\{A(\theta)\Lambda^{-1}(\theta)[F - \eta(\theta, V)] - A(\theta)\dot{V}\right\}
\end{align}

\textbf{最终结果}：
\begin{equation}
\lambda = (A\Lambda^{-1}A^T)^{-1}(A\Lambda^{-1}(F - \eta) - A\dot{V}) \tag{11.58}
\end{equation}

其中为了简洁，我们省略了$(\theta)$和$(V)$的依赖关系。

\subsection{约束力的计算}

\textbf{机器人施加到约束上的wrench}：
\begin{equation}
F_{\text{tip}} = A^T(\theta)\lambda \tag{11.59}
\end{equation}

\textbf{物理意义}：
\begin{itemize}
\item $F_{\text{tip}}$是机器人施加到环境上的力
\item 这个力与约束方向相关
\item 由拉格朗日乘数$\lambda$决定
\end{itemize}

\section{投影矩阵的推导}

\subsection{问题：如何分离运动子空间和力子空间？}

\textbf{目标}：
\begin{itemize}
\item 将$n$维力空间划分为两个子空间
\item 运动控制子空间：维度$n-k$（自由运动方向）
\item 力控制子空间：维度$k$（约束方向）
\end{itemize}

\textbf{方法}：使用投影矩阵

\subsection{投影矩阵的定义}

\textbf{投影矩阵$P$}：
\begin{equation}
P = I - A^T(A\Lambda^{-1}A^T)^{-1}A\Lambda^{-1} \tag{11.60}
\end{equation}

\textbf{性质}：
\begin{itemize}
\item $P$的秩为$n-k$（假设$n=6$，则秩为$6-k$）
\item $P$将任意wrench $F$投影到运动控制子空间
\item $I - P$将任意wrench $F$投影到力控制子空间
\end{itemize}

\subsection{投影矩阵的推导}

\textbf{目标}：从约束动力学中消除约束力项

\textbf{步骤1：从约束动力学求解$\lambda$}

从方程(11.56)：
\begin{equation}
F = \Lambda(\theta)\dot{V} + \eta(\theta, V) + A^T(\theta)\lambda
\end{equation}

我们已经知道：
\begin{equation}
\lambda = (A\Lambda^{-1}A^T)^{-1}(A\Lambda^{-1}(F - \eta) - A\dot{V})
\end{equation}

\textbf{步骤2：代入约束动力学}

将$\lambda$代入：
\begin{align}
F &= \Lambda\dot{V} + \eta + A^T\lambda \\
  &= \Lambda\dot{V} + \eta + A^T(A\Lambda^{-1}A^T)^{-1}(A\Lambda^{-1}(F - \eta) - A\dot{V})
\end{align}

\textbf{步骤3：展开}

展开$A^T\lambda$项：
\begin{align}
A^T\lambda &= A^T(A\Lambda^{-1}A^T)^{-1}A\Lambda^{-1}(F - \eta) - A^T(A\Lambda^{-1}A^T)^{-1}A\dot{V}
\end{align}

代入：
\begin{align}
F &= \Lambda\dot{V} + \eta + A^T(A\Lambda^{-1}A^T)^{-1}A\Lambda^{-1}(F - \eta) - A^T(A\Lambda^{-1}A^T)^{-1}A\dot{V}
\end{align}

\textbf{步骤4：整理}

将关于$F$的项移到左边：
\begin{align}
F - A^T(A\Lambda^{-1}A^T)^{-1}A\Lambda^{-1}F &= \Lambda\dot{V} + \eta - A^T(A\Lambda^{-1}A^T)^{-1}A\Lambda^{-1}\eta - A^T(A\Lambda^{-1}A^T)^{-1}A\dot{V}
\end{align}

左边：
\begin{align}
[I - A^T(A\Lambda^{-1}A^T)^{-1}A\Lambda^{-1}]F &= PF
\end{align}

右边：
\begin{align}
&\Lambda\dot{V} + \eta - A^T(A\Lambda^{-1}A^T)^{-1}A\Lambda^{-1}\eta - A^T(A\Lambda^{-1}A^T)^{-1}A\dot{V} \\
&= [I - A^T(A\Lambda^{-1}A^T)^{-1}A\Lambda^{-1}](\Lambda\dot{V} + \eta) \\
&= P(\Lambda\dot{V} + \eta)
\end{align}

\textbf{最终结果}：
\begin{equation}
P(\theta)F = P(\theta)(\Lambda(\theta)\dot{V} + \eta(\theta, V)) \tag{11.61}
\end{equation}

\textbf{物理意义}：
\begin{itemize}
\item 左边：投影后的力（运动控制子空间）
\item 右边：投影后的动力学（运动控制子空间）
\item 约束力项$A^T\lambda$被投影消除了
\end{itemize}

\subsection{投影矩阵的性质}

\textbf{性质1：$P$是投影矩阵}

验证$P^2 = P$：
\begin{align}
P^2 &= [I - A^T(A\Lambda^{-1}A^T)^{-1}A\Lambda^{-1}]^2 \\
    &= I - 2A^T(A\Lambda^{-1}A^T)^{-1}A\Lambda^{-1} + A^T(A\Lambda^{-1}A^T)^{-1}A\Lambda^{-1}A^T(A\Lambda^{-1}A^T)^{-1}A\Lambda^{-1} \\
    &= I - 2A^T(A\Lambda^{-1}A^T)^{-1}A\Lambda^{-1} + A^T(A\Lambda^{-1}A^T)^{-1}A\Lambda^{-1} \\
    &= I - A^T(A\Lambda^{-1}A^T)^{-1}A\Lambda^{-1} = P
\end{align}

因此$P$是投影矩阵。

\textbf{性质2：$P$和$I-P$正交}

验证$P(I-P) = 0$：
\begin{align}
P(I-P) &= P - P^2 = P - P = 0
\end{align}

因此$P$和$I-P$是正交的投影矩阵。

\textbf{性质3：$P$的秩}

\begin{itemize}
\item $A(\theta) \in \mathbb{R}^{k \times n}$，秩为$k$
\item $A^T(\theta) \in \mathbb{R}^{n \times k}$，秩为$k$
\item $A\Lambda^{-1}A^T \in \mathbb{R}^{k \times k}$，秩为$k$（假设约束独立）
\item $A^T(A\Lambda^{-1}A^T)^{-1}A\Lambda^{-1}$的秩为$k$
\item 因此$P = I - A^T(A\Lambda^{-1}A^T)^{-1}A\Lambda^{-1}$的秩为$n-k$
\end{itemize}

\section{混合控制器的构建}

\subsection{基本思想}

\textbf{目标}：
\begin{itemize}
\item 在运动控制子空间（自由方向）使用运动控制器
\item 在力控制子空间（约束方向）使用力控制器
\item 将两者组合起来
\end{itemize}

\subsection{运动控制器}

\textbf{任务空间运动控制器}（从计算力矩控制律(11.47)导出）：
\begin{equation}
F_{\text{motion}} = \tilde{\Lambda}(\theta)\frac{d}{dt}([Ad_{X^{-1}X_d}]V_d) + K_p X_e + K_i\int_0^t X_e(t)dt + K_d V_e
\end{equation}

其中：
\begin{itemize}
\item $X_d$：期望的末端执行器位姿
\item $V_d$：期望的末端执行器速度
\item $X_e = \log(X^{-1}X_d)$：位置误差
\item $V_e$：速度误差
\item $K_p, K_i, K_d$：控制增益矩阵
\end{itemize}

\textbf{投影到运动子空间}：
\begin{equation}
F_{\text{motion,proj}} = P(\theta)F_{\text{motion}}
\end{equation}

\subsection{力控制器}

\textbf{任务空间力控制器}（公式(11.51)）：
\begin{equation}
F_{\text{force}} = F_d + K_{fp}F_e + K_{fi}\int_0^t F_e(t)dt
\end{equation}

其中：
\begin{itemize}
\item $F_d$：期望的力
\item $F_e = F_d - F_{\text{tip}}$：力误差
\item $K_{fp}, K_{fi}$：控制增益矩阵
\end{itemize}

\textbf{投影到力子空间}：
\begin{equation}
F_{\text{force,proj}} = (I - P(\theta))F_{\text{force}}
\end{equation}

\subsection{组合控制器}

\textbf{完整的混合控制器}：
\begin{equation}
F = F_{\text{motion,proj}} + F_{\text{force,proj}} + \text{补偿项}
\end{equation}

展开：
\begin{align}
F &= P(\theta)F_{\text{motion}} + (I - P(\theta))F_{\text{force}} + \tilde{\eta}(\theta, V_b) \\
  &= P(\theta)\left[\tilde{\Lambda}(\theta)\frac{d}{dt}([Ad_{X^{-1}X_d}]V_d) + K_p X_e + K_i\int_0^t X_e(t)dt + K_d V_e\right] \\
  &\quad + (I - P(\theta))\left[F_d + K_{fp}F_e + K_{fi}\int_0^t F_e(t)dt\right] + \tilde{\eta}(\theta, V_b)
\end{align}

\textbf{转换到关节空间}：

使用虚功原理：$\tau = J_b^T(\theta)F$

\begin{equation}
\tau = J_b^T(\theta)\Bigg[P(\theta)\left(\tilde{\Lambda}(\theta)\frac{d}{dt}([Ad_{X^{-1}X_d}]V_d) + K_p X_e + K_i\int_0^t X_e(t)dt + K_d V_e\right) + (I - P(\theta))\left(F_d + K_{fp}F_e + K_{fi}\int_0^t F_e(t)dt\right) + \tilde{\eta}(\theta, V_b)\Bigg] \tag{11.61}
\end{equation}

\section{控制器的解耦性}

\subsection{为什么两个控制器是解耦的？}

\textbf{关键}：投影矩阵$P$和$I-P$是正交的

\textbf{证明}：
\begin{itemize}
\item $P(I-P) = 0$（已证明）
\item 因此$P$和$I-P$投影到正交的子空间
\item 运动控制器和力控制器作用在不同的子空间上
\item 因此它们是解耦的
\end{itemize}

\textbf{结果}：
\begin{itemize}
\item 运动控制器的误差动力学和稳定性分析独立于力控制器
\item 力控制器的误差动力学和稳定性分析独立于运动控制器
\item 混合控制器继承了各自子空间上单个控制器的性质
\end{itemize}

\section{总结}

\subsection{推导流程}

\begin{enumerate}
\item \textbf{先验知识}：
\begin{itemize}
\item 任务空间动力学：$F = \Lambda(\theta)\dot{V} + \eta(\theta, V)$
\item Pfaffian约束：$A(\theta)V = 0$
\end{itemize}

\item \textbf{约束动力学}：
\begin{itemize}
\item 使用拉格朗日乘数法：$F = \Lambda(\theta)\dot{V} + \eta(\theta, V) + A^T(\theta)\lambda$
\item 约束的时间导数：$A(\theta)\dot{V} + \dot{A}(\theta)V = 0$
\end{itemize}

\item \textbf{拉格朗日乘数}：
\begin{itemize}
\item 求解$\lambda = (A\Lambda^{-1}A^T)^{-1}(A\Lambda^{-1}(F - \eta) - A\dot{V})$
\end{itemize}

\item \textbf{投影矩阵}：
\begin{itemize}
\item 定义$P = I - A^T(A\Lambda^{-1}A^T)^{-1}A\Lambda^{-1}$
\item 投影后的动力学：$PF = P(\Lambda\dot{V} + \eta)$
\end{itemize}

\item \textbf{混合控制器}：
\begin{itemize}
\item 运动控制器投影到运动子空间：$PF_{\text{motion}}$
\item 力控制器投影到力子空间：$(I-P)F_{\text{force}}$
\item 组合：$\tau = J_b^T[PF_{\text{motion}} + (I-P)F_{\text{force}} + \eta]$
\end{itemize}
\end{enumerate}

\subsection{关键公式}

\begin{enumerate}
\item \textbf{约束动力学}：
\begin{equation}
F = \Lambda(\theta)\dot{V} + \eta(\theta, V) + A^T(\theta)\lambda
\end{equation}

\item \textbf{拉格朗日乘数}：
\begin{equation}
\lambda = (A\Lambda^{-1}A^T)^{-1}(A\Lambda^{-1}(F - \eta) - A\dot{V})
\end{equation}

\item \textbf{投影矩阵}：
\begin{equation}
P = I - A^T(A\Lambda^{-1}A^T)^{-1}A\Lambda^{-1}
\end{equation}

\item \textbf{混合控制器}：
\begin{equation}
\tau = J_b^T(\theta)\Bigg[P(\theta)F_{\text{motion}} + (I - P(\theta))F_{\text{force}} + \tilde{\eta}(\theta, V_b)\Bigg]
\end{equation}
\end{enumerate}

\subsection{物理意义}

\begin{itemize}
\item \textbf{约束动力学}：考虑环境约束对系统动力学的影响
\item \textbf{拉格朗日乘数}：表示约束力的大小
\item \textbf{投影矩阵}：分离运动子空间和力子空间
\item \textbf{混合控制器}：在不同方向上分别使用运动控制和力控制
\end{itemize}

\section{通俗解释}

\subsection{核心思想}

\textbf{类比}：想象你在推一个箱子：

\begin{itemize}
\item \textbf{水平方向}（自由运动）：你可以控制箱子的位置和速度
\item \textbf{垂直方向}（有约束）：箱子被地面支撑，你不能控制垂直位置，但可以控制垂直力（推多用力）
\end{itemize}

\textbf{混合控制}：
\begin{itemize}
\item 在水平方向：使用运动控制（控制位置和速度）
\item 在垂直方向：使用力控制（控制力的大小）
\item 两者同时进行，互不干扰
\end{itemize}

\subsection{为什么需要投影？}

\textbf{问题}：如何确保运动控制器只影响自由方向，力控制器只影响约束方向？

\textbf{答案}：使用投影矩阵！

\begin{itemize}
\item \textbf{投影矩阵$P$}：将力投影到自由运动方向
\item \textbf{投影矩阵$I-P$}：将力投影到约束方向
\item 两者正交，互不干扰
\end{itemize}

\textbf{结果}：
\begin{itemize}
\item 运动控制器产生的力被投影到自由方向
\item 力控制器产生的力被投影到约束方向
\item 两者可以同时工作，不会冲突
\end{itemize}

\section{详细例子：擦黑板}

\subsection{场景设置}

\textbf{任务}：机器人用橡皮擦擦黑板

\textbf{约束}：
\begin{itemize}
\item 不能倾斜：$\omega_x = 0, \omega_y = 0$
\item 不能穿透：$v_z = 0$
\item 总共$k = 3$个约束
\end{itemize}

\textbf{约束矩阵}：
\begin{equation}
A(\theta) = \begin{bmatrix}
1 & 0 & 0 & 0 & 0 & 0 \\
0 & 1 & 0 & 0 & 0 & 0 \\
0 & 0 & 0 & 0 & 0 & 1
\end{bmatrix}
\end{equation}

\subsection{控制目标}

\textbf{运动控制}（在自由方向）：
\begin{itemize}
\item $v_x = k_1$（水平移动速度）
\item $v_y = 0$（不横向移动）
\item $\omega_z = 0$（不旋转）
\end{itemize}

\textbf{力控制}（在约束方向）：
\begin{itemize}
\item $f_z = k_2 < 0$（保持与黑板的接触力）
\item $m_x = 0, m_y = 0$（保持水平）
\end{itemize}

\subsection{控制器实现}

\textbf{步骤1：计算投影矩阵}

\begin{equation}
P = I - A^T(A\Lambda^{-1}A^T)^{-1}A\Lambda^{-1}
\end{equation}

\textbf{步骤2：设计运动控制器}

在运动子空间（$PF_{\text{motion}}$）：
\begin{itemize}
\item 控制$v_x, v_y, \omega_z$
\item 使用PID控制器跟踪期望轨迹
\end{itemize}

\textbf{步骤3：设计力控制器}

在力子空间（$(I-P)F_{\text{force}}$）：
\begin{itemize}
\item 控制$f_z, m_x, m_y$
\item 使用PI控制器跟踪期望力
\end{itemize}

\textbf{步骤4：组合}

\begin{equation}
\tau = J_b^T[PF_{\text{motion}} + (I-P)F_{\text{force}} + \eta]
\end{equation}

\section{实现注意事项}

\subsection{约束矩阵的确定}

\textbf{问题}：如何知道约束矩阵$A(\theta)$？

\textbf{方法}：
\begin{enumerate}
\item \textbf{环境建模}：根据环境特性确定约束
\item \textbf{实时估计}：使用力反馈识别约束方向
\item \textbf{低增益控制}：使用低反馈增益，使控制器对约束不确定性更容忍
\end{enumerate}

\subsection{奇异点问题}

\textbf{问题}：在奇异点，雅可比矩阵$J_b(\theta)$不可逆

\textbf{影响}：
\begin{itemize}
\item 任务空间质量矩阵$\Lambda(\theta)$变得很大
\item 投影矩阵计算可能不稳定
\end{itemize}

\textbf{解决方法}：
\begin{itemize}
\item 避免奇异配置
\item 使用伪逆$J_b^{\dagger}$代替$J_b^{-1}$
\item 添加阻尼项
\end{itemize}

\subsection{被动柔顺性}

\textbf{问题}：关节和连杆中的柔顺性不可避免

\textbf{影响}：
\begin{itemize}
\item 实际约束不是完全刚性的
\item 约束矩阵$A(\theta)$可能不准确
\end{itemize}

\textbf{解决方法}：
\begin{itemize}
\item 在机器人结构中构建被动柔顺性
\item 使用低反馈增益
\item 实时估计约束方向
\end{itemize}

\section{数学细节补充}

\subsection{约束时间导数的详细推导}

\textbf{约束}：
\begin{equation}
A(\theta)V = 0
\end{equation}

\textbf{对时间求导}：
\begin{align}
\frac{d}{dt}[A(\theta)V] &= \frac{d}{dt}[A(\theta)]V + A(\theta)\frac{dV}{dt} \\
&= \dot{A}(\theta)V + A(\theta)\dot{V} = 0
\end{align}

\textbf{关键点}：
\begin{itemize}
\item $\dot{A}(\theta)$是约束矩阵对时间的导数
\item 由于$A(\theta)$依赖于$\theta$，而$\theta$依赖于时间，所以$\dot{A}(\theta) = \frac{\partial A}{\partial \theta}\dot{\theta}$
\item 这反映了约束可能随着机器人配置的变化而变化
\end{itemize}

\subsection{投影矩阵的几何意义}

\textbf{力空间的分解}：

$n$维力空间可以分解为两个正交子空间：
\begin{itemize}
\item \textbf{运动控制子空间}：维度$n-k$，由$P$投影
\item \textbf{力控制子空间}：维度$k$，由$I-P$投影
\end{itemize}

\textbf{几何解释}：
\begin{itemize}
\item 运动控制子空间是约束子空间的正交补空间
\item 力控制子空间是约束子空间本身
\item 两个子空间正交，因此投影是解耦的
\end{itemize}

\subsection{拉格朗日乘数的物理意义}

\textbf{拉格朗日乘数$\lambda$}：
\begin{equation}
\lambda = (A\Lambda^{-1}A^T)^{-1}(A\Lambda^{-1}(F - \eta) - A\dot{V})
\end{equation}

\textbf{物理意义}：
\begin{itemize}
\item $\lambda$表示约束力的大小
\item 约束力$F_{\text{tip}} = A^T\lambda$是机器人施加到环境上的力
\item $\lambda$的大小取决于：
\begin{itemize}
\item 外部力$F$
\item 科里奥利项$\eta$
\item 加速度$\dot{V}$
\end{itemize}
\end{itemize}

\section{总结}

\subsection{核心概念}

\begin{enumerate}
\item \textbf{约束动力学}：考虑环境约束对系统动力学的影响
\item \textbf{拉格朗日乘数}：表示约束力的大小
\item \textbf{投影矩阵}：分离运动子空间和力子空间
\item \textbf{混合控制器}：在不同方向上分别使用运动控制和力控制
\end{enumerate}

\subsection{关键公式总结}

\begin{enumerate}
\item \textbf{约束动力学}：
\begin{equation}
F = \Lambda(\theta)\dot{V} + \eta(\theta, V) + A^T(\theta)\lambda
\end{equation}

\item \textbf{约束的时间导数}：
\begin{equation}
A(\theta)\dot{V} + \dot{A}(\theta)V = 0
\end{equation}

\item \textbf{拉格朗日乘数}：
\begin{equation}
\lambda = (A\Lambda^{-1}A^T)^{-1}(A\Lambda^{-1}(F - \eta) - A\dot{V})
\end{equation}

\item \textbf{投影矩阵}：
\begin{equation}
P = I - A^T(A\Lambda^{-1}A^T)^{-1}A\Lambda^{-1}
\end{equation}

\item \textbf{投影后的动力学}：
\begin{equation}
PF = P(\Lambda\dot{V} + \eta)
\end{equation}

\item \textbf{混合控制器}：
\begin{equation}
\tau = J_b^T(\theta)\Bigg[P(\theta)F_{\text{motion}} + (I - P(\theta))F_{\text{force}} + \tilde{\eta}(\theta, V_b)\Bigg]
\end{equation}
\end{enumerate}

\subsection{实际应用}

混合运动-力控制器适用于：
\begin{itemize}
\item 与环境接触的任务（如装配、打磨、写字）
\item 需要在不同方向上分别控制运动和力的任务
\item 有明确约束方向的任务
\end{itemize}

\subsection{设计步骤}

\begin{enumerate}
\item \textbf{识别约束}：确定环境约束矩阵$A(\theta)$
\item \textbf{计算投影}：计算投影矩阵$P$和$I-P$
\item \textbf{设计控制器}：
\begin{itemize}
\item 在运动子空间设计运动控制器
\item 在力子空间设计力控制器
\end{itemize}
\item \textbf{组合实现}：将两个控制器组合，并转换到关节空间
\item \textbf{补偿项}：添加科里奥利和重力补偿
\end{enumerate}

\end{document}